\section{Nghiên cứu người dùng}

\subsection{Đối tượng người dùng mục tiêu}
% Trình bày:
% - Người dùng mục tiêu (Chủ nuôi thú cưng bận rộn theo Proposal)
% - Đặc điểm nhân khẩu học và thói quen chăm sóc thú cưng
% - Ngữ cảnh sử dụng thực tế (Context of use)

\subsection{Phương pháp nghiên cứu}
% Trình bày phương pháp thực hiện (Phỏng vấn sâu, Khảo sát, Quan sát...):
% - Mục đích
% - Số lượng và tiêu chí người tham gia
% - Quy trình thực hiện
% - Dữ liệu thu thập

\subsection{Kết quả nghiên cứu}
% Trình bày các phát hiện và dữ liệu thực tế:
% - Người dùng thường gặp khó khăn khi...
% - Người dùng kỳ vọng...
% - Người dùng nhầm lẫn hoặc lo lắng về...

\subsection{Nhu cầu và điểm đau của người dùng}
% Tổng hợp nhu cầu và điểm đau chính của người dùng (kế thừa từ Proposal):
% - Điểm đau: Chờ xác nhận lâu, thất lạc dặn dò đặc biệt/dị ứng, thiếu minh bạch tiến độ...
% - Nhu cầu: Đặt lịch tức thì, tự động lưu hồ sơ dị ứng/thuốc, theo dõi tiến độ thời gian thực...

\subsection{Chân dung người dùng}
% Trình bày Persona dựa trên dữ liệu phỏng vấn thực tế và Proposal:
% - Primary Persona: Chị Lan (28 tuổi, Nhân viên văn phòng, nuôi chó Poodle có tiền sử dị ứng sữa tắm)
% - Mục tiêu, Hành vi, Điểm đau, Nhu cầu

